% !TeX root = ../papers/nc-p3.tex
\section{第三長方形の全正則性と無印の四元群}
\label{sec:third-rectangle-regularity}
\begin{theorem}\label{thm:third-rectangle-all-regular}
第三長方形の台88,89,90の全点が標準Koszul左右AS正則である。
各始点は四つの中心四次式上の自由階数256の加群となる。
従って定理\ref{thm:third-rectangle-orbits}の三次元商、その周期層・反対階数層・台の特殊化は
全て正則代数の添字保持同型類を与える。
\end{theorem}
\begin{proof}
$a=A,h=AD,b=x_1$ とし、$ahb(h-1)(b-1)$ を逆にする。この三変数環上で経路を長さ5まで計算し、
文字0,1の四次中心核を求める。局所化の単元主元だけで自由列数3,2を得る。
文字0の最初と最後の自由列、文字1の両自由列の四式を選ぶ。
\begin{enumerate}
\item 始点2では、台の表示をそのまま使うと四式の商の長さ13までの計算にも追加の逆元が一つも要らない。
従って全点で $(1+z+z^2+z^3)^4$ の次元列を得る。二位置移動 $a\mapsto h/a$ により始点0も全点で従う。
\item 始点1では文字を位置順に $(2,2,0,0)$ だけ四元群で平行移動し、同じ自由列選択法で四式を取り直す。
中心商の計算に要る新しい既約因子は
\[
 F=abh+ab-6bh+4h
\]
だけである。指定始点を変えない同型 $(a,b)\mapsto(-a,b/(b-1))$ でこの表示を移すと、
必要な非零条件は単元倍を除いて $F'=-abh-ab-2bh-4h$ になる。
$F+F'=-8bh$ は可逆なので両方が失敗する点はない。始点3も二位置移動で従う。
\end{enumerate}
\path{computations/verify_third_rectangle_regularity.py} が二つの表示の素因子集合を空集合と $\{F\}$ にまで再計算する。
一般点から元の代数には全次数の標準Hilbert下界がある。各始点で成功する中心式を一般に選び直し、
四始点の有限個の開条件を同時に満たす四式を取れる。中心商による自由加群の上界との比較で自由性を得る。
元のテンソルの階数と二つの中心正則元から、補題\ref{lem:parallel-residue-border}と同じ
左右Koszul列の完全性が従う。
\end{proof}

\begin{proposition}\label{prop:rectangular-intrinsic-klein}
四文字空間に射影的に作用して窓の直線を保つ群は、第三長方形の全点、および
台145の両水平面の $a=y=-1$ 以外の全点で、標準の四元群そのものである。
ここには両面の台143の曲線と台142への交わりを含む。
特にこれらの窓に別の四元群の印を選んで異なる文字多重度の枝に移すことはできない。
\end{proposition}
\begin{proof}
既に復元した四枠により各位置で単項行列である。第三長方形で恒等以外の置換は正規表示を $a\mapsto-a$ に移すので、
指定始点で同じ係数を保てない。水平面でも一般の台で追加の置換が係数を保つには同時逆数変換の固定点が必要である。
一零条件の同じ向きの台を保つ置換は恒等だけである。二零条件では $y^2-y+1$ で割った体で、
台の置換のうち係数比の全対角指標が1となるものを数えると各符号とも恒等一つになる。

対角群は台の位置別重み行列から決まる。台88,89,90,142,143,145のいずれも非零Smith因子は
$1^{11},2,2$ である。各位置のスカラーを除けば有限部の位数は4になり、
四元群の二つの独立な文字がこの全てを実現する。これは射影文字空間上の主張であり、
この群に忠実に作用する別の四元群も同じ部分群に一致する。
よってその固有文字線は共通で、印の変更では群の自己同型と各位置の文字平行移動だけが残る。
\path{computations/verify_rectangle_klein_stabilizers.py} が台ごとの有限部と周期1の二境界での全置換を検算する。
\end{proof}
